\paragraph{\trEquivariant{}情形。}\label{sub:equivariant-case}
我们从一些简单的几何准备开始。设 $M$ 是一个带 $\ZZ/2$-作用 $\iota_M$ 的\trExactSymplectic\trManifold。$M$ 中一个\textbf{不变的带\trFrame{}的\trExact\trLagrange\trSphere}（invariant framed exact Lagrangian sphere）是指一个 $\iota_M$-不变的\trExact\trLagrange\trSubmanifold{}$V$，连同一种\trEquivariant\trDiffeomorphism{}$v\colon S\to V$，其中 $S$ 是某个\trInnerProduct{}空间中的单位\trSphere，而该空间配备有一个\trLinear{}的 $\ZZ/2$-作用。由于\trLagrange\trTubularNeighborhood{}定理可以做得与有限群作用\trCompatible，并且\trDehnTwist{}的局部模型关于 $O_{n+1}$ 是\trEquivariant{}的，沿这种\trSphere{}的\trDehnTwist{}$\tau_L$ 自然是一个\trEquivariant{}的\trExactSymplectic\trAutomorphism。给定一个带有\trFibrewise\trInvolution{}的\trExact\trLefschetzFibration，与之相伴的诸\trParallelTransport{}映射都是\trEquivariant{}的，诸\trVanishingCycle{}都是\trEquivariant{}带\trFrame{}的\trSphere，而且 \trPicard--\trLefschetz{}公式 \eqref{eq:symplectic-picard-lefschetz} 的一个适当的\trEquivariant{}版本成立。反过来，给定 $M$ 中一个或几个\trEquivariant{}带\trFrame{}的\trSphere，则可以安排得使\cref{thm:single-vanishing-cycle-fibration}与\cref{thm:realizing-distinguished-basis}的那些\trLefschetzFibration{}带有\trFibrewise\trInvolution。
现在假设我们处于 \cref{sub:equivariant-fukaya-category} 的情形之中，于是\trEquivariant\trFukayaCategory{}便可以定义。我们把它记作 $\scr{A}=\scr{F}(M)^{\mathrm{naive}}$，把它上面的 $\ZZ/2$-作用记作 $\epsilon_{\scr{A}}$。设有一个\trEquivariant{}带\trFrame{}的\trExact\trLagrange\trSphere{}$V=L_0$，它可以被做成一个\trEquivariant\trBrane{}$L^{\#}_0=(L_0,\alpha^{\#}_0,P^{\#}_0,f_0)$。我们还想要另外两个\trEquivariant\trBrane{}$L^{\#}_k=(L_k,\alpha^{\#}_k,P^{\#}_k,f_k)$（$k=1,2$），连同一条穿过诸不变\trExact\trLagrange\trSubmanifold{}的\trIsotopy{}$\tau_V(L_1)\simeq L_2$，并且可以把这条\trIsotopy{}做得与诸\trEquivariant\trBraneStructure\trCompatible。在这种情形下，我们可以考虑由诸 $L^{\#}_k$ 给出的 $\epsilon_{\scr{A}}$-不动\trObject{}$Y_k\in\Ob\scr{A}$，以及与之相伴的\trEquivariant{}\trYoneda{}\trModule{}$\scr{Y}_k\in\Ob\mathit{mod}^{\mathrm{eq}}(\scr{A})$。正如 \cref{sub:equivariant-a-infinity-modules} 中所指出的，\trTwisted{} \trAinfty-模 $\scr{T}_{Y_0}(\scr{Y}_1)$ 仍是\trEquivariant{}的。\cref{thm:twist-tau-isomorphism-modules}的\trEquivariant{}类似物说，$(c,k)$ 给出诸\trEquivariant{}\trModule{}之间的一个\trMorphism
\begin{equation}\label{eq:equivariant-module-morphism-teq}
[t^{\mathrm{eq}}]\in\Hom_{H^0(\mathit{mod}^{\mathrm{eq}}(\scr{A}))}\bigl(\scr{T}_{Y_0}(\scr{Y}_1),\scr{Y}_2\bigr),
\end{equation}
它是该\trCategory{}中的一个\trIsomorphism。尽管这里没有任何新东西，一段简短的解释可能仍然有帮助。首先，$t^{\mathrm{eq}}$ 由与\cref{thm:quasi-representative-criterion}(i) 中相同的那些公式定义。现在回忆，若忘掉 $\ZZ/2$-作用，$\scr{A}$ 就变得与 $\scr{F}(M)$ \trQuasiequivalent。更精确地说，通过把每个\trLagrange\trBrane{}做成预\trEquivariant\trBrane（这总是可能的，见\cref{thm:equivariant-brane-remark}），我们可以把 $\scr{F}(M)$ 等同于 $\scr{A}$ 的一个\trFullSubcategory{}——它与整体\trQuasiequivalent；而且由此得到的\trFunctor
\begin{equation}\label{eq:forget-equivariance-functor}
\mathit{mod}^{\mathrm{eq}}(\scr{A})\longrightarrow\mathit{mod}(\scr{A})\longrightarrow\mathit{mod}(\scr{F}(M))
\end{equation}
是\trCohomology\trFullyFaithful{}的（第一部分，即忘掉\trEquivariant{}性，在链的层面上\trFull{}且\trFaithful；第二部分，即限制到 $\scr{F}(M)\subset\scr{A}$ 上，据\cref{thm:restriction-quasi-equivalence}是一个\trQuasiequivalence）。此外，还可以安排\trEmbedding{}$\scr{F}(M)\subset\scr{F}(M)^{\mathrm{naive}}$，使得诸\trLagrange\trBrane{}$L^{\#}_k$ 被映到它们\trEquivariant{}的对应物。于是，$t^{\mathrm{eq}}$ 在 \eqref{eq:forget-equivariance-functor} 之下的像就是先前考虑过的那个\trModule{}\trHomomorphism{}$t$，这就表明 $[t^{\mathrm{eq}}]$ 是一个\trIsomorphism，正如上面所陈述的那样。
\begin{lemma}\label{thm:teq-invariant-part}
$[t^{\mathrm{eq}}]$ 属于 \eqref{eq:equivariant-module-morphism-teq} 中那个 $\Hom$ 群的 $\ZZ/2$-不变部分。
\end{lemma}
人们可以相当有理由地宣称这是平凡的：我们先前已经证明 $[t]$ 是\trCanonical{}的，因此它在我们几何设定中固有的一切对称性之下都应当不变，而同样的论述也适用于它的\trEquivariant{}版本。我们仍然想要看一看细节，以确证那里没有潜伏着不受欢迎的意外。可以安排得使诸\trLefschetzFibration{}$E^c\to S^c$ 与 $E^{k,r}\to S^{k,r}$，连同它们的\trBoundaryCondition{}一起，都带有与 $M$ 上的那个作用\trCompatible{}的\trFibrewise{} $\ZZ/2$-作用。所涉及的全部\trLagrange\trSubmanifold{}都是\trEquivariant\trBrane，由 $\scr{F}(M)^{\mathrm{naive}}$ 的构造，它们自带 $\ZZ/2$-不变的\trFloerDatum。剩下的选取是\trRelative\trPerturbationDatum{}的选取。在这一点上，我们宁愿避开\trEquivariant\trTransversality{}问题，转而如下论证。设 $L^{\#}_1$、$L^{\#}_2$ 更一般地是预\trEquivariant\trBrane。首先考虑带有其\trFibrewise\trInvolution{}的 $E^c\to S^c$。把用来定义 $c$ 的那个\trLagrange\trBoundaryCondition{}记作 $F$，把它的\trRelative\trBraneStructure{}记作 $F^{\#}$，并把相关的\trRelative\trPerturbationDatum{}记作 $(K,J)$。我们可以把\trFibrewise\trInvolution{}施加于 $(K,J)$ 与 $F^{\#}$ 这两者之上，然后把后者的 \trPinStructure{}用\trLineBundle{}$\lambda^{\mathrm{top}}(TF^v)$ 作张量。在这个修改过的情形中对\trInhomogeneous\trPseudoholomorphic\trSection{}进行计数所得的结果，是一个类 $CF^0\bigl(\iota_M(L^{\#}_1\otimes\lambda^{\mathrm{top}}(TL_1)),\iota_M(L^{\#}_2\otimes\lambda^{\mathrm{top}}(TL_2))\bigr)$，按定义它恰恰就是 $\epsilon_{\scr{A}}(c)$。现在，在 $L^{\#}_1$ 与 $L^{\#}_2$ 都是\trEquivariant\trBrane{}的情形，诸\trIsomorphism{}$f_k$ 给出一个自然的等同
\begin{equation*}
CF^0\bigl(\iota_M(L^{\#}_1\otimes\lambda^{\mathrm{top}}(TL_1)),\iota_M(L^{\#}_2\otimes\lambda^{\mathrm{top}}(TL_2))\bigr)=CF^0(L^{\#}_1,L^{\#}_2)
\end{equation*}
它来自诸 \trPinStructure{}的\trIsomorphism{}$P^{\#}_k\to\iota_M^*\bigl(P^{\#}_k\otimes\lambda^{\mathrm{top}}(TL_k)\bigr)$（$k=1,2$）。在我们的情形，这些\trIsomorphism{}与\trIsotopy{}$\tau_V(L_1)\simeq L_2$ \trCompatible，这意味着它们给出了\trLagrange\trBoundaryCondition{}$F$ 的\trRelative{} \trPinStructure{}之间的一个\trIsomorphism。记住了这一点，显然 $\epsilon_{\scr{A}}(c)$ 与 $c$ 的差别仅仅在于\trRelative\trPerturbationDatum{}的选取。$k$ 的情形完全平行。于是，\eqref{eq:cocycle-homotopy-independence} 的构造适用（插入适当的符号），给出 $(\gamma,\kappa)$ 使得
\begin{equation*}
\begin{gathered}
\mu^1_{\scr{A}}(\gamma)+c-\epsilon_{\scr{A}}(c)=0,\\
(-1)^{|c_0|-1}\bigl(\mu^1_{\scr{A}}\bigl(\kappa(c_0)\bigr)-\kappa\bigl(\mu^1_{\scr{A}}(c_0)\bigr)+\mu^2_{\scr{A}}(\gamma,c_0)\bigr)\\
\quad-\epsilon_{\scr{A}}\bigl(\kappa(\epsilon_{\scr{A}}(c_0))\bigr)+\kappa(c_0)=0
\quad\text{对一切 }c_0\in CF^*(L^{\#}_0,L^{\#}_1).
\end{gathered}
\end{equation*}
有了这些，\cref{thm:quasi-representative-criterion}(ii) 的论证便产生一个介于 $t^{\mathrm{eq}}$ 与其在 $\ZZ/2$-作用之下的像之间的\trHomotopy，从而证明在\trCohomology{}的层面上，$[t^{\mathrm{eq}}]$ 确实是不变的。
